\section{Side-Effect Analysis}\label{sec:side_effect}

Given a command, it is useful to know
which already existing lists might be affected during the execution of the command.
Previous work has defined precise and complex static analyses for that
(see for instance the aliasing and side-effect analyses in~\cite{SpotoMP10}).
Here, for simplicity and lack of space,
we present a simpler and cheaper side-effect analysis that, nevertheless,
is often precise enough on simple programs such as smart contracts.

Possible side-effects are an abstract interpretation of execution traces.
Therefore, the side-effect analysis is presented as a bottom-up denotational
semantics that uses an interpretation\ie inductive information about
the side-effect of the method calls triggered by the execution of a command.
We do not discuss this further here, since it is standard technology
in semantics~\cite{Winskel93} and abstract interpretation~\cite{CousotC77}.
An example of computation is reported later.
%
\begin{definition}[Side-Effects Analysis]\label{def:side_effect}
  The set of element types of the already existing lists possibly modified during the
  execution of a command $\com$ is overapproximated by $\sideeffect{\com}$,
  defined in Fig.~\ref{fig:side_effect} by using an \emph{interpretation} $I$ that
  provides the possible side-effects of the methods.
  %
  \begin{figure}[t]
  \begin{align*}
    \sideeffect{\assign{\mathit{v}}{\new{C}{v^*}}}&=I(C.C)\\
    \sideeffect{\assign{\mathit{v}}{\exp}}&=\varnothing\text{ for any other $\exp$}\\
    \sideeffect{\push{v_1}{v_2}}&=\{t\mid v_1\text{ has type }\<List<>t\<\text{>}>\}\\
    \sideeffect{\call{v}{m}{v^*}}&=\cup_{j=1}^k I(C_j.m)\\
    \noalign{\centering\text{where $C_1.m\cdots C_k.m$ are the methods possibly called here}}
    \sideeffect{\ife{v}{v}{\com_1}{\com_2}}&=\sideeffect{\com_1}\cup\sideeffect{\com_2}\\
    \sideeffect{\for{v}{v}{\com}}&=\sideeffect{\com}\\
    \sideeffect{\sequence{\com_1}{\com_2}}&=\sideeffect{\com_1}\cup\sideeffect{\com_2}\\
    \sideeffect{\com}&=\varnothing\text{, for all other commands.}
  \end{align*}
  \caption{The side-effect rules. They use an interpretation $I$ for the methods.}
  \label{fig:side_effect}
  \end{figure}
\end{definition}
%
We recall that a method call in object-oriented languages may actually call a subset
$\{C_1.m\cdots C_k.m\}$ of the redefinitions of its static target method.
A sound overapproximation of $\{C_1.m\cdots C_k.m\}$ can always be easily inferred statically,
as the set of \emph{all} possible redefinitions. More precise analyses exist~\cite{TipP00} as well,
but they are probably overkill for minimalistic code such as smart contracts,
that very rarely redefine methods.

The computation of the side-effect analysis
starts from the empty interpretation $I_0$, that provides $\varnothing$ as
side-effects of every method. The bodies of the methods
(hence also of the constructors) of the program are then analyzed with the rules
in Fig.~\ref{fig:side_effect}, using $I=I_0$, to get an improved interpretation $I_1$.
This process is repeated to yield $I_2$ and so on, until a fixpoint $\intsideeffect$
is reached. This is guaranteed to happen after a finite number of iterations
since the rules in Fig.~\ref{fig:side_effect} are monotonic and the set of element types
of lists used in each given program is finite.

\begin{example}\label{ex:side_effect}
  Fig.~\ref{fig:side_effect_examples} reports $\intsideeffect$ for the smart contracts from
  Sec.~\ref{sec:examples}\ie the element types of the lists potentially modified during
  the execution of each of their constructors or methods. Note that these approximations are correct
  but sometimes imprecise. For instance, $\intsideeffect(\<Ponzi1.Ponzi1>)=\{\<Contract>\}$
  is correct but could be improved to a more precise $\intsideeffect(\<Ponzi1.Ponzi1>)=\varnothing$,
  since no existing list of \<Contract> is modified by the execution of that constructor:
  the list modified at line~\ref{ponzi1:creator} of Fig.~\ref{fig:ponzi1}
  is a fresh new list created at the previous line, not a list that existed before the
  execution of the constructor. These imprecisions could be overcome with more complex analyses,
  but the results in Fig.~\ref{fig:side_effect_examples} are good enough for our purposes.
  %
  \begin{figure}[t]
    \begin{center}
    \begin{subfigure}{0.42\textwidth}
      \begin{align*}
        \intsideeffect(\<Ponzi1.Ponzi1>)&=\{\<Contract>\}\\
        \intsideeffect(\<Ponzi1.invest>)&=\{\<Contract>\}
      \end{align*}
      \caption{The side-effect analysis of the smart contract in Fig.~\ref{fig:ponzi1}
      (similar analysis also for that in Fig.~\ref{fig:ponzi3}).}
    \end{subfigure}
    \qquad
    \begin{subfigure}{0.42\textwidth}
      \begin{align*}
        \intsideeffect(\<Ponzi2.Ponzi2>)&=\{\<Contract>\}\\
        \intsideeffect(\<Ponzi2.invest>)&=\{\<Contract>\}\\
        \intsideeffect(\<Ponzi2.withdraw>)&=\varnothing
      \end{align*}
      \caption{The side-effect analysis of the smart contract in Fig.~\ref{fig:ponzi2}.}
    \end{subfigure}\\
    \begin{subfigure}{\textwidth}
      \begin{align*}
        \intsideeffect(\<TicketBatchSale.TicketBatchSale>)&=\varnothing\\
        \intsideeffect(\<Round.Round>)&=\varnothing\\
        \intsideeffect(\<Lottery1.Lottery1>)&=\{\<Round>\}\\
        \intsideeffect(\<Lottery1.buy>)&=\{\<Round,TicketBatchSale>\}\\
        \intsideeffect(\<Lottery1.drawWinner>)&=\varnothing
      \end{align*}
      \caption{The side-effect analysis of the smart contract in Fig.~\ref{fig:lottery1}
      (the analysis remains similar after the modification in Fig.~\ref{fig:lottery2}).}
    \end{subfigure}
    \end{center}
    \caption{The side-effect analysis of the smart contracts from Sec.~\ref{sec:examples}.}
    \label{fig:side_effect_examples}
  \end{figure}
\end{example}

In the following, we just assume that $I_\mathsf{se}$ has been computed.
